%%═══════════════════════════════════════════════════════════════════
%% Lecture 01 — Introduction: Special Relativity as the Foundation
%%              of Electrodynamics; Galilean Transformations and
%%              Their Failure for Maxwell's Equations
%% 77401 — Analytical Electrodynamics
%% Date:
%% Lecturer: Prof. Michael Smolkin
%%═══════════════════════════════════════════════════════════════════

\section{Lecture 01 — Introduction: Galilean Transformations and the
  Failure of Non-Relativistic Symmetry}\label{sec:lec01}

\begin{center}
\begin{tabular}{ll}
  \textbf{Date:}\quad 22/03/2026   & \\
\end{tabular}
\end{center}
\hrule\vspace{1.5em}

%%──────────────────────────────────────────────────────────────────
\subsection*{Overview}

This lecture opens the course by motivating Analytical Electrodynamics
as the field-theoretic analogue of Analytical Mechanics: where the
latter reformulates finite-degree-of-freedom Newtonian dynamics via
the action principle (Lagrange, Hamilton), this course does the same
for fields — systems with \emph{infinitely many} degrees of freedom.
The central goal is to derive Maxwell's equations from an action
principle, identify the symmetries of the theory via Noether's theorem,
and ultimately understand why special relativity is not merely
compatible with electrodynamics but is, in a precise sense, demanded
by it.  The lecture then begins the technical programme by reviewing
the structure of Galilean spacetime, proving that the Galilean group is
indeed a group, and demonstrating that Maxwell's equations are
\emph{not} invariant under Galilean transformations — a contradiction
that forces us to abandon the Galilean picture in favour of Lorentzian
spacetime and the Lorentz group.

%%──────────────────────────────────────────────────────────────────
\subsection*{Definitions}

\dfn{Inertial reference frame}{A reference frame in which a body subject to no external forces moves
with constant velocity (including zero).  Equivalently, Newton's
second law holds in its standard form in such a frame.  The existence
of at least one inertial frame is taken as an axiom; all frames moving
with constant velocity relative to a known inertial frame are also
inertial.}

\dfn{Homogeneity of spacetime}{The laws of physics are invariant under arbitrary translations in both
space and time: no spatial location and no instant of time is
privileged.}

\dfn{Isotropy of space}{The laws of physics are invariant under arbitrary rotations of the
spatial axes.  There is no preferred direction in space.  (Isotropy
applies to space; homogeneity applies to spacetime.)}

\dfn{Principle of Relativity}{The \emph{equations of motion} — i.e., the laws of physics — take
identical form in every inertial reference frame.  All inertial frames
are equivalent; no preferred inertial frame exists.}

\dfn{Galilean transformation}{Given two inertial frames $S$ and $S'$ with $S'$ moving at constant
velocity $\mathbf{u}$ relative to $S$, with origins coinciding at
$t = 0$, the Galilean transformation relating an event
$(\mathbf{r}, t)$ in $S$ to its coordinates $(\mathbf{r}', t')$ in
$S'$ is
\begin{align}
  \mathbf{r}' &= \mathbf{r} - \mathbf{u}\,t, \label{eq:lec01:GT_r}\\
  t'           &= t.                           \label{eq:lec01:GT_t}
\end{align}
Time is absolute (frame-independent) in the Galilean picture.}

\dfn{Group}{A set $G$ together with a binary operation $\circ$ is a
\emph{group} if it satisfies:
\begin{enumerate}
  \item \textbf{Closure.} For all $g_1, g_2 \in G$,\;
        $g_1 \circ g_2 \in G$.
  \item \textbf{Identity.} There exists $e \in G$ such that
        $e \circ g = g \circ e = g$ for all $g \in G$.
  \item \textbf{Inverse.} For every $g \in G$ there exists
        $g^{-1} \in G$ such that $g \circ g^{-1} = g^{-1} \circ g = e$.
  \item \textbf{Associativity.} For all $g_1, g_2, g_3 \in G$,\;
        $(g_1 \circ g_2) \circ g_3 = g_1 \circ (g_2 \circ g_3)$.
\end{enumerate}}

%%──────────────────────────────────────────────────────────────────
\subsection*{Key Results}

\thm{Galilean transformations form a group}{The set of all Galilean boosts
$\{G_{\mathbf{u}} : \mathbf{u} \in \mathbb{R}^3\}$, under the
operation of composition, forms a group — the \emph{Galilean group}
(boost subgroup).  The composition law is
\begin{equation}
  G_{\mathbf{u}''} = G_{\mathbf{u}'} \circ G_{\mathbf{u}},
  \qquad \mathbf{u}'' = \mathbf{u} + \mathbf{u}'.
  \label{eq:lec01:galilean_group_law}
\end{equation}}

\thm{Newton's second law is Galilean-invariant}{Under a Galilean transformation \eqref{eq:lec01:GT_r}–\eqref{eq:lec01:GT_t},
the equation $m\,\ddot{\mathbf{r}} = -\nabla V(\mathbf{r})$ retains
identical form:
\begin{equation}
  m\,\frac{d^2\mathbf{r}'}{dt'^2} = -\nabla' V(\mathbf{r}').
  \label{eq:lec01:newton_invariant}
\end{equation}}

\thm{Galilean field-transformation rules (from Lorentz force invariance)}{Demanding that the Lorentz force law
$m\ddot{\mathbf{r}} = q\mathbf{E} + \frac{q}{c}\mathbf{v}\times\mathbf{B}$
be invariant under Galilean transformations (working in Gaussian/CGS
units) forces the electromagnetic fields to transform as
\begin{align}
  \mathbf{B}' &= \mathbf{B},                                   \label{eq:lec01:Btransform}\\
  \mathbf{E}' &= \mathbf{E} + \frac{1}{c}\,\mathbf{u}\times\mathbf{B}. \label{eq:lec01:Etransform}
\end{align}}

\thm{Maxwell's equations are \emph{not} Galilean-invariant}{The field-transformation rules
\eqref{eq:lec01:Btransform}–\eqref{eq:lec01:Etransform}, which follow
from demanding Galilean invariance of the Lorentz force, are
inconsistent with Gauss's law $\nabla'\cdot\mathbf{E}' = 0$ (in a
source-free region).  Specifically, computing
$\nabla'\cdot\mathbf{E}'$ yields a term proportional to
$\mathbf{u}\cdot\partial_t\mathbf{E}$, which is generically
non-zero.  Therefore the Galilean transformation law for
spacetime is physically incorrect, and must be replaced by the
Lorentz transformation.}

The single most important equation of the lecture is the
incompatibility result:
\begin{equation}
  \boxed{
    \nabla' \cdot \mathbf{E}' \;=\;
    \frac{1}{c}\,\mathbf{u}\cdot\frac{\partial\mathbf{E}}{\partial t}
    \;\neq\; 0 \quad \text{in general},
  }
  \label{eq:lec01:inconsistency}
\end{equation}
which signals that Galilean spacetime is incompatible with Maxwell's
theory.

%%──────────────────────────────────────────────────────────────────
\subsection*{Derivations}

%% DERIVATION 1: Galilean group closure
\subsubsection*{Galilean boosts form a group: closure}

Let $G_{\mathbf{u}}$ denote the boost by velocity $\mathbf{u}$:
\begin{align}
  \mathbf{r}' &= \mathbf{r} - \mathbf{u}\,t, &
  t'           &= t.
  \label{eq:lec01:boost-closure-start}
\end{align}
Let $G_{\mathbf{u}'}$ act next on $(\mathbf{r}', t')$:
\begin{align}
  \mathbf{r}'' &= \mathbf{r}' - \mathbf{u}'\,t'
               = \bigl(\mathbf{r} - \mathbf{u}\,t\bigr)
                 - \mathbf{u}'\,t
               = \mathbf{r} - (\mathbf{u}+\mathbf{u}')\,t,\\
  t''          &= t' = t.
\end{align}
The result has the same form as a single Galilean boost with velocity
$\mathbf{u}'' = \mathbf{u} + \mathbf{u}'$, proving closure.
The identity element is $G_{\mathbf{0}}$ ($\mathbf{r}' = \mathbf{r}$,
$t' = t$); the inverse of $G_{\mathbf{u}}$ is $G_{-\mathbf{u}}$;
associativity follows from function composition.

%% DERIVATION 2: Galilean invariance of Newton's second law
\subsubsection*{Newton's second law is Galilean-invariant}

Consider a particle of mass $m$ moving in a scalar potential $V$:
\begin{equation}
  m\,\frac{d^2\mathbf{r}}{dt^2} = -\nabla_{\mathbf{r}} V(\mathbf{r}).
  \label{eq:lec01:newt}
\end{equation}

\noindent\textbf{Left-hand side.}
Since $t' = t$, time derivatives are unchanged.  From
$\mathbf{r}' = \mathbf{r} - \mathbf{u} t$,
\begin{equation}
  \frac{d\mathbf{r}'}{dt'} = \frac{d\mathbf{r}}{dt} - \mathbf{u},
  \label{eq:lec01:velocity-transform}
\end{equation}
the velocity shifts by $-\mathbf{u}$ (Galilean velocity addition).
Taking a second derivative (and recalling $\mathbf{u}$ is constant):
\begin{equation}
  \frac{d^2\mathbf{r}'}{dt'^2} = \frac{d^2\mathbf{r}}{dt^2}.
  \label{eq:lec01:accel_invariant}
\end{equation}
Acceleration is the same in both frames.

\noindent\textbf{Right-hand side.}
The chain rule gives
\begin{equation}
  \frac{\partial}{\partial x} V(\mathbf{r})
  = \frac{\partial x'}{\partial x}\frac{\partial V}{\partial x'}
  + \frac{\partial y'}{\partial x}\frac{\partial V}{\partial y'}
  + \frac{\partial z'}{\partial x}\frac{\partial V}{\partial z'}.
  \label{eq:lec01:chain-rule-potential}
\end{equation}
From $x' = x - u_x t$, $y' = y - u_y t$, $z' = z - u_z t$, one
finds $\partial x'/\partial x = 1$, $\partial y'/\partial x = 0$,
$\partial z'/\partial x = 0$.  Hence
$\partial V/\partial x = \partial V/\partial x'$, and likewise
for $y$ and $z$, so
\begin{equation}
  \nabla_{\mathbf{r}} V(\mathbf{r}) = \nabla_{\mathbf{r}'} V(\mathbf{r}').
  \label{eq:lec01:gradient-invariant}
\end{equation}
Equations \eqref{eq:lec01:accel_invariant} and the above together show
that \eqref{eq:lec01:newt} becomes \eqref{eq:lec01:newton_invariant}
in $S'$, with no change of form.

%% DERIVATION 3: Field transformation laws from Lorentz force invariance
\subsubsection*{Galilean field transformation from Lorentz force invariance}

The Lorentz force in $S$ (CGS units) is
\begin{equation}
  m\,\ddot{\mathbf{r}} = q\mathbf{E} + \frac{q}{c}\,\mathbf{v}\times\mathbf{B},
  \label{eq:lec01:lorentz_S}
\end{equation}
where $\mathbf{E}$ and $\mathbf{B}$ are external fields.
By the principle of relativity under Galilean transformations the same
equation must hold in $S'$:
\begin{equation}
  m\,\ddot{\mathbf{r}}' = q\mathbf{E}' + \frac{q}{c}\,\mathbf{v}'\times\mathbf{B}'.
  \label{eq:lec01:lorentz_Sp}
\end{equation}
Since $\ddot{\mathbf{r}}' = \ddot{\mathbf{r}}$ (shown above), the
right-hand sides of \eqref{eq:lec01:lorentz_S} and
\eqref{eq:lec01:lorentz_Sp} must be equal.
The Galilean velocity transformation gives
$\mathbf{v}' = \mathbf{v} - \mathbf{u}$.  Substituting:
\begin{equation}
  q\mathbf{E}' + \frac{q}{c}(\mathbf{v}-\mathbf{u})\times\mathbf{B}'
  = q\mathbf{E} + \frac{q}{c}\,\mathbf{v}\times\mathbf{B}.
  \label{eq:lec01:lorentz-force-match}
\end{equation}
Rearranging, collecting terms with and without $\mathbf{v}$:
\begin{equation}
  q\mathbf{E}' - \frac{q}{c}\,\mathbf{u}\times\mathbf{B}'
  + \frac{q}{c}\,\mathbf{v}\times(\mathbf{B}'-\mathbf{B})
  = q\mathbf{E}.
  \label{eq:lec01:rearranged}
\end{equation}
The fields $\mathbf{E}$, $\mathbf{B}$, $\mathbf{E}'$, $\mathbf{B}'$
are generated by \emph{external} sources and are independent of
$\mathbf{v}$ (the velocity of the test particle).  Because $\mathbf{v}$
is freely controllable, the coefficient of $\mathbf{v}$ in
\eqref{eq:lec01:rearranged} must vanish identically:
\begin{equation}
  \mathbf{B}' - \mathbf{B} = \mathbf{0}
  \implies
  \mathbf{B}' = \mathbf{B}.
  \label{eq:lec01:Btransform-derived}
\end{equation}
Substituting this back into \eqref{eq:lec01:rearranged}:
\begin{equation}
  \mathbf{E}' = \mathbf{E} + \frac{1}{c}\,\mathbf{u}\times\mathbf{B}.
  \label{eq:lec01:Etransform-derived}
\end{equation}
These are equations \eqref{eq:lec01:Btransform}–\eqref{eq:lec01:Etransform}.

%% DERIVATION 4: Inconsistency with Gauss's law
\subsubsection*{Galilean field transforms are inconsistent with Maxwell's equations}

Work in a source-free region: $\nabla\cdot\mathbf{E} = 0$ in $S$, and
by the principle of relativity $\nabla'\cdot\mathbf{E}' = 0$ must also
hold in $S'$.  Compute $\nabla'\cdot\mathbf{E}'$ using the derived
transformation:
\begin{equation}
  \nabla'\cdot\mathbf{E}' = \nabla'\cdot\mathbf{E}
                           + \frac{1}{c}\,\nabla'\cdot(\mathbf{u}\times\mathbf{B}).
  \label{eq:lec01:gauss-primed-expand}
\end{equation}
Under a Galilean transformation $\partial/\partial x'^{\,i} =
\partial/\partial x^i$ (the Jacobian is the identity), so
$\nabla'\cdot\mathbf{E} = \nabla\cdot\mathbf{E} = 0$.
For the second term, using the vector identity
$\nabla\cdot(\mathbf{A}\times\mathbf{C})
 = \mathbf{C}\cdot(\nabla\times\mathbf{A}) - \mathbf{A}\cdot(\nabla\times\mathbf{C})$
with $\mathbf{u}$ constant:
\begin{equation}
  \nabla\cdot(\mathbf{u}\times\mathbf{B})
  = -\mathbf{u}\cdot(\nabla\times\mathbf{B}).
  \label{eq:lec01:div-cross-identity}
\end{equation}
The Ampère–Maxwell law (in CGS, source-free) gives
$\nabla\times\mathbf{B} \propto \partial\mathbf{E}/\partial t$,
so:
\begin{equation}
  \nabla'\cdot\mathbf{E}'
  = -\frac{1}{c}\,\mathbf{u}\cdot\frac{\partial\mathbf{E}}{\partial t}.
  \label{eq:lec01:gauss-contradiction}
\end{equation}
Because $\partial\mathbf{E}/\partial t \neq \mathbf{0}$ in general (e.g.\
in an electromagnetic wave), and $\mathbf{u}$ can be chosen with any
direction, this is generically non-zero — in contradiction with the
requirement $\nabla'\cdot\mathbf{E}' = 0$.  The Galilean law of
transformation of inertial frames is therefore physically untenable:
it must be replaced by the Lorentz transformation.

%%──────────────────────────────────────────────────────────────────
\subsection*{Examples}

\ex{Closure of Galilean boosts: explicit computation}{Let $G_{\mathbf{u}}$ boost from $S$ to $S'$ and $G_{\mathbf{u}'}$
boost from $S'$ to $S''$.  With $S'$ moving at $\mathbf{u}$ and
$S''$ moving at $\mathbf{u}'$ relative to $S'$:
\begin{align}
  \mathbf{r}'' &= \mathbf{r}' - \mathbf{u}' t'
               = \mathbf{r} - \mathbf{u} t - \mathbf{u}' t
               = \mathbf{r} - (\mathbf{u}+\mathbf{u}')\,t,
               \quad t'' = t.
\end{align}
The composition equals a single boost $G_{\mathbf{u}+\mathbf{u}'}$,
confirming the group law $\mathbf{u}'' = \mathbf{u} + \mathbf{u}'$
and group closure.}

\ex{Integers under addition as a group}{Smolkin used $(\mathbb{Z}, +)$ as a concrete illustration of the group
axioms: the identity is $0$, and the inverse of any integer $n$ is
$-n$.  Closure ($\mathbb{Z}$ is closed under $+$) and associativity
are immediate.  This parallels the Galilean boosts, where $0$
corresponds to the trivial boost and $-\mathbf{u}$ to its inverse.}

%%──────────────────────────────────────────────────────────────────
\subsection*{Connections}

\begin{itemize}
  \item \textbf{Analytical Mechanics.}  Smolkin explicitly frames this
    course as the field-theory analogue of Analytical Mechanics: the
    action principle and Noether's theorem for fields generalize what
    Lagrange and Hamilton did for finite-degree-of-freedom systems.
    Prior exposure to the Lagrangian formulation (including the
    Euler–Lagrange equations) is assumed.
  \item \textbf{Quantum Field Theory.}  Smolkin noted that the
    formalism developed here — classical field theory with an action
    principle — is the direct precursor to QFT, which combines
    special relativity with quantum mechanics and underpins all of
    modern elementary-particle physics.
  \item \textbf{Special Relativity (from Mechanics course).}
    Students are expected to know the Lorentz transformation from
    their mechanics course.  The next lectures will review and extend
    this knowledge; Lorentz transformations will replace Galilean ones
    as the correct symmetry of spacetime.
  \item \textbf{Landau \& Lifshitz, Vol.\ 2 (\emph{Classical Theory of
    Fields}).}  Smolkin stated that the majority of the course follows
    this text closely.  The title ``Classical Theory of Fields'' reflects
    that, when the book was written, there were exactly two classical
    field theories: electrodynamics and gravity.  The course covers only
    the electrodynamics part.
  \item \textbf{Jackson, \emph{Classical Electrodynamics}.}
    Recommended as an equally authoritative alternative reference,
    covering most topics in this course.
\end{itemize}

%%──────────────────────────────────────────────────────────────────
\subsection*{To Remember}

\begin{itemize}
  \item The Galilean transformation asserts absolute time ($t' = t$)
    and the position rule $\mathbf{r}' = \mathbf{r} - \mathbf{u}t$;
    it implies Galilean velocity addition
    $\mathbf{v}' = \mathbf{v} - \mathbf{u}$.
  \item Acceleration is an \emph{absolute} quantity under Galilean
    transformations: $\ddot{\mathbf{r}}' = \ddot{\mathbf{r}}$.
    This makes Newton's second law (with velocity-independent forces)
    Galilean-invariant.
  \item The Galilean boosts $\{G_{\mathbf{u}}\}$ form a group under
    composition, with composition law
    $G_{\mathbf{u}'}\circ G_{\mathbf{u}} = G_{\mathbf{u}+\mathbf{u}'}$.
  \item Every symmetry of physics has the mathematical structure of a
    group; the symmetry group of non-relativistic mechanics is the
    Galilean group.
  \item Demanding Galilean invariance of the Lorentz force law forces
    $\mathbf{B}' = \mathbf{B}$ and
    $\mathbf{E}' = \mathbf{E} + \mathbf{u}\times\mathbf{B}/c$.
  \item These field-transformation rules are inconsistent with
    $\nabla\cdot\mathbf{E}=0$ in $S'$ (source-free Gauss's law) because
    $\nabla'\cdot(\mathbf{u}\times\mathbf{B}) = -\mathbf{u}\cdot
    (\nabla\times\mathbf{B}) \propto \mathbf{u}\cdot\partial_t\mathbf{E}
    \neq 0$ in general.
  \item \textbf{Conclusion:} Maxwell's equations are not
    Galilean-invariant; the correct transformation is the Lorentz
    transformation, to be studied in the next lectures.
  \item The course works exclusively in \textbf{Gaussian (CGS) units}
    because relativistic factors of $c$ appear symmetrically; in SI
    units they are obscured.
\end{itemize}

%%──────────────────────────────────────────────────────────────────
\subsection*{Exam / HW Hints}

\begin{itemize}
  \item Smolkin strongly recommended treating homework seriously and
    working through it independently (or with understanding), as it
    directly prepares for the exam.
  \item He noted that AI tools and solutions from peers are
    permissible but warned explicitly: you must read and understand
    whatever you submit; copying without understanding is both
    ineffective for the exam and a poor learning strategy.
  \item Office hours are by appointment via email
    (\texttt{michael.smolkin@mail.huji.ac.il}); demand peaks in the
    week before the exam, so he advised scheduling early.
  \item Assessment: final exam 90\%; HW bonus up to 15\% (maximum
    grade 100\%); submitting $\geq 70\%$ of assignments is required
    for the full bonus.
\end{itemize}
